\section{Introduction}
Scientific writing assistants are increasingly used to improve draft quality, but many tools still operate as one-shot editors with limited control over iterative refinement.
This paper introduces a modular reviewer pipeline that alternates between issue detection, targeted edits, and quality scoring over multiple rounds.
Compared with plain prompt-based rewriting, our method keeps explicit history and supports early stopping when no further improvement is observed.
The core research question is whether lightweight orchestration can produce stable quality gains with low implementation complexity.

\section{Related Work}
Prior work on automated writing support includes grammar correction systems, style transfer models, and LLM-based revision agents.
Grammar correction methods often optimize token-level fluency, while revision agents target discourse-level coherence and argument structure.
Recent graph-based orchestration frameworks have shown strong flexibility in multi-step reasoning tasks.
However, few studies provide practical engineering baselines that are easy to run locally and inspect end-to-end.

\section{Method}
Our framework is implemented as a directed graph with five main roles: initializer, reviewer, editor, critic, and aggregator.
Given input LaTeX text, the initializer splits content by section markers and stores structured state.
For each section, the reviewer emits issue objects containing fields such as section id, severity, and problem description.
The editor applies deterministic edits based on these issues, and the critic outputs a scalar score in $[0,1]$.
The aggregator updates current text, tracks best text, and records a complete modification history.

\subsection{State Definition}
The state includes immutable input, mutable current text, per-section slices, issue list, running score, best score, and stop counters.
Formally, we maintain
\[
S_t = \{x_{orig}, x_t, \mathcal{Q}_t, h_t, b_t, i_t\},
\]
where $x_t$ is current draft, $\mathcal{Q}_t$ is issue set, $h_t$ is edit history, and $b_t$ stores the best-so-far artifact.

\subsection{Routing Policy}
Section routing follows:
\[
\text{if } idx < N \Rightarrow \text{reviewer}, \quad \text{else} \Rightarrow \text{iteration\_step}.
\]
Iteration routing follows:
\[
\text{if } iter \geq I_{max} \text{ or } no\_improve \geq K \Rightarrow \text{end}, \text{ else reviewer}.
\]

\section{Experimental Setup}
We run all experiments in a local conda environment with fixed dependency versions to ensure reproducibility.
Each run processes the same synthetic manuscript and reports iteration count, best score, and number of accepted edits.
To control randomness in evaluation, we fix random seeds when needed and isolate deterministic checks in unit tests.

\section{Results}
Table~\ref{tab:main} summarizes representative outcomes across three runs.
\begin{table}[h]
\centering
\begin{tabular}{lccc}
\hline
Run & Iterations & Best Score & Accepted Edits \\
\hline
A & 2 & 0.89 & 6 \\
B & 2 & 0.91 & 7 \\
C & 3 & 0.93 & 8 \\
\hline
\end{tabular}
\caption{Example iterative optimization outcomes.}
\label{tab:main}
\end{table}

Qualitatively, revisions improve local clarity and reduce repetitive phrasing.
Failure cases remain when issues require domain knowledge beyond sentence-level rewriting.

\section{Ablation Study}
We remove one component at a time to estimate contribution:
\begin{itemize}
\item Without critic: no reliable acceptance signal, quality fluctuates.
\item Without history: debugging and rollback are difficult.
\item Without section loop: edits become coarse and less controllable.
\end{itemize}
The full pipeline consistently provides the best trade-off between quality and interpretability.

\section{Discussion}
The framework is intentionally simple and should be viewed as an engineering baseline rather than a final product.
Its modular design enables future upgrades, such as replacing mock reviewer outputs with LLM calls and adding semantic consistency checks.
A practical limitation is that section splitting currently relies on a narrow LaTeX pattern and does not yet support all structural variants.

\section{Conclusion}
We presented a lightweight graph-based reviewer that supports iterative, section-aware text refinement.
Despite its simplicity, the system demonstrates stable behavior, traceable decisions, and low setup overhead.
Future work will include stronger scoring functions, richer edit operators, and benchmark evaluation on real manuscripts.
